\documentclass[11pt]{article}
\usepackage[margin=1in]{geometry}
\usepackage{amsmath,amssymb}
\usepackage{booktabs}
\usepackage{graphicx}
\usepackage{url}
\usepackage{hyperref}
\usepackage{xcolor}

\title{An Empirical Comparison of Knowledge-Distillation Objectives\\for Compressing a Pretrained Language Model}
\author{Tate Staples\thanks{Experiments, implementation, and drafting assisted by Claude (Anthropic).}}
\date{July 2026}

\begin{document}
\maketitle

\begin{abstract}
Knowledge distillation (KD) transfers the behavior of a large ``teacher''
language model into a smaller ``student.'' The literature offers many
objectives---forward and reverse Kullback--Leibler divergence on output
logits, symmetric divergences such as the generalized Jensen--Shannon
divergence and total variation distance, intermediate-representation
matching, and sequence-level distillation on teacher-generated text---but
few controlled comparisons at small scale. We distill the pretrained
GPT-2 (124M parameters) into a 29M-parameter, 4-layer student under an
equal step budget for seven objectives, evaluating WikiText-2 perplexity,
LAMBADA final-word accuracy, and HellaSwag and ARC-Easy multiple-choice
accuracy. We additionally investigate an auxiliary
representation-alignment target motivated by a user-supplied reference to
an Anthropic ``J-Space'' paper; we document that no such paper exists in
the public record and instead evaluate the closest defensible
interpretation. \textbf{[RESULTS SUMMARY TO BE FILLED FROM
runs/summary.json]}
\end{abstract}

\section{Introduction}
Large language models are expensive to serve; distillation
\cite{hinton2015distilling} compresses them by training a small student
to imitate a large teacher rather than only the raw data. Since Hinton
et al.'s original formulation, the field has produced a family of
objectives that differ chiefly in (i) \emph{what} is matched (output
distributions, intermediate states, or generated sequences) and (ii)
\emph{which divergence} measures the mismatch. This paper implements the
most common objectives in a single controlled harness and compares them
when distilling a pretrained GPT-2 teacher on a fixed compute budget.

\section{Methods}
\label{sec:methods}
Let $p(\cdot \mid x)$ be the teacher's next-token distribution and
$q_\theta(\cdot \mid x)$ the student's. All students share one
architecture, initialization, optimizer, and step budget; only the loss
differs.

\paragraph{Hard-label baseline (\texttt{hard}).}
Standard next-token cross-entropy on the data,
$\mathcal{L}_{\mathrm{CE}} = -\log q_\theta(y_t \mid x_{<t})$,
with no teacher involvement.

\paragraph{Forward-KL logit distillation (\texttt{kd}).}
The classic objective \cite{hinton2015distilling}: with temperature $T$
and mixing weight $\alpha$,
\begin{equation}
\mathcal{L} = \alpha\,\mathcal{L}_{\mathrm{CE}}
 + (1-\alpha)\,T^2\,\mathrm{KL}\!\left(p^{1/T} \,\|\, q_\theta^{1/T}\right),
\end{equation}
where $p^{1/T}$ denotes the temperature-flattened distribution. Forward
KL is \emph{mass-covering}: the student is penalized wherever the teacher
assigns probability the student misses.

\paragraph{Reverse-KL distillation (\texttt{rkl}).}
$\mathrm{KL}(q_\theta \,\|\, p)$ is \emph{mode-seeking} and was argued by
MiniLLM \cite{gu2024minillm} to suit generative students that cannot
cover every teacher mode. We use the teacher-forced variant (computed on
data prefixes rather than on-policy student rollouts) to keep the
comparison budget-matched.

\paragraph{Generalized Jensen--Shannon divergence (\texttt{jsd}).}
Following GKD \cite{agarwal2024gkd},
$\mathrm{JSD}_\beta(p, q) = \beta\,\mathrm{KL}(p\|m) +
(1-\beta)\,\mathrm{KL}(q\|m)$ with $m = \beta p + (1-\beta) q$;
$\beta{=}0.5$ gives the symmetric JSD, and the family interpolates
between forward ($\beta\!\to\!1$) and reverse ($\beta\!\to\!0$) KL.

\paragraph{Total variation distance (\texttt{tvd}).}
$\tfrac12 \sum_v |p_v - q_v|$, a bounded symmetric alternative studied in
f-divergence distillation \cite{wen2023fdivergence}.

\paragraph{Hidden-state matching (\texttt{hidden}).}
TinyBERT/DistilBERT-style \cite{jiao2020tinybert,sanh2019distilbert}
intermediate-layer transfer: each student layer $s$ is mapped to teacher
layer $t(s)$ (uniform spacing), a learned linear projection lifts the
student width (384) to the teacher width (768), and a scale-normalized
MSE between projected student and teacher states is added to the
forward-KL objective.

\paragraph{Sequence-level KD (\texttt{seqkd}).}
Kim \& Rush \cite{kim2016sequence}: the teacher generates full sequences
(here, sampled completions of 16-token WikiText-2 prompts,
top-$k{=}50$), and the student trains with plain cross-entropy on this
synthetic corpus.

\paragraph{J-space auxiliary target (\texttt{jspace}).}
\textbf{[TO BE WRITTEN AFTER RESEARCH: what was found about the claimed
Anthropic paper, the chosen interpretation, and its formula.]}

\section{Experimental Setup}
\label{sec:setup}
\textbf{Teacher.} GPT-2 124M (12 layers, 768-dim), pretrained weights in
HF format.\footnote{Weights obtained from the legacy Hugging Face S3
mirror; the experiment environment blocks huggingface.co.}
\textbf{Student.} 4 layers, 384-dim, 6 heads, GPT-2 BPE vocabulary
(50{,}257), $\sim$29M parameters ($\sim$23\% of teacher), randomly
initialized, identical seed across methods.
\textbf{Data.} WikiText-2 (word-level release), packed into 128-token
blocks; 4{,}000 training blocks.
\textbf{Optimization.} AdamW, lr $3\times10^{-4}$ with linear warmup
(5\%) and linear decay, weight decay 0.01, gradient clipping 1.0,
batch size 8, 600 steps ($\approx$0.6M training tokens), $T{=}2$,
$\alpha{=}0.5$. Single seed (compute-constrained; CPU-only).
\textbf{Evaluation.} WikiText-2 test perplexity on 200 held-out blocks;
LAMBADA final-word greedy accuracy (300 examples); HellaSwag and
ARC-Easy accuracy via length-normalized log-likelihood over answer
choices (300 examples each).

\section{Results}
\label{sec:results}
\textbf{[TABLE + ANALYSIS TO BE FILLED FROM runs/summary.json]}

\section{Discussion and Limitations}
\textbf{[TO BE WRITTEN: interpretation, capacity-gap notes, small-budget
caveats, single-seed caveat, teacher-forced rkl caveat.]}

\begin{thebibliography}{10}
\bibitem{hinton2015distilling} G.~Hinton, O.~Vinyals, J.~Dean.
Distilling the knowledge in a neural network. \emph{arXiv:1503.02531}, 2015.
\bibitem{kim2016sequence} Y.~Kim, A.~Rush. Sequence-level knowledge
distillation. \emph{EMNLP}, 2016.
\bibitem{sanh2019distilbert} V.~Sanh, L.~Debut, J.~Chaumond, T.~Wolf.
DistilBERT, a distilled version of BERT. \emph{arXiv:1910.01108}, 2019.
\bibitem{jiao2020tinybert} X.~Jiao et al. TinyBERT: Distilling BERT for
natural language understanding. \emph{Findings of EMNLP}, 2020.
\bibitem{gu2024minillm} Y.~Gu, L.~Dong, F.~Wei, M.~Huang. MiniLLM:
Knowledge distillation of large language models. \emph{ICLR}, 2024.
\bibitem{agarwal2024gkd} R.~Agarwal et al. On-policy distillation of
language models: Learning from self-generated mistakes. \emph{ICLR}, 2024.
\bibitem{wen2023fdivergence} Y.~Wen, Z.~Li, W.~Du, L.~Mou.
f-divergence minimization for sequence-level knowledge distillation.
\emph{ACL}, 2023.
\end{thebibliography}

\end{document}
